\chapter{数据与方法设计}

整体方法划分为数据采集、文本处理、主题分析和标准化解释四层。
采集层维护来源白名单、种子地址和访问日志；文本处理层完成正文抽取、字段规范化、精确去重、近重复识别和中文分词；
主题分析层输出词频、TF--IDF、共词关系和候选主题；标准化解释层识别对象、过程、接口、数据、服务、检测评价或管理活动。

\section{合规采集与可追溯记录}

数据源按政府部门、标准平台、研究机构和行业组织分类登记。
每个文档至少保存来源、URL、发布时间、采集时间、正文摘要和内容指纹，避免把转载次数误写为独立需求强度。

\section{文本处理与主题分析}

\subsection{特征表示}

设文档集合为 $D$，词项 $t$ 在文档 $d$ 中的 TF--IDF 权重写为

\begin{equation}
  w_{t,d}=\operatorname{tf}(t,d)\log\frac{|D|}{1+\operatorname{df}(t)}.
  \label{eq:tfidf}
\end{equation}

\subsubsection{数据不足时的回退策略}

当语料数量和文本质量不足时，流程回退到“关键词统计、共词分析与人工编码”，不强行使用 LDA 或 K-Means。

\begin{figure}[htbp]
  \centering
  \includegraphics[width=0.92\textwidth]{research-route.png}
  \caption{低空经济标准需求研究技术路线}
  \label{fig:research-route}
\end{figure}

图\ref{fig:research-route}用于测试外部 PNG、中文题注和交叉引用的 PDF/DOCX 转换。
